\begin{proposition}[2.8.5]
\label{IV.2.8.5-fr}
Siendo las hipótesis sobre $X$ e $Y$ las de \textup{(2.8.1)}, sea $Z'$ un
subpreesquema cerrado de $X_\eta$. Existe entonces un único subpreesquema
cerrado $\overline{Z'}$ de $X$ que sea $Y$-\textup{plano} y tal que
$i^{-1}(\overline{Z'})=Z'$.
\end{proposition}

\begin{proof}
Si $\mathcal J'$ es el ideal cuasicoherente de $\mathcal O_{X_\eta}$
que define $Z'$, basta aplicar (2.8.1) al caso en que
$\mathcal F=\mathcal O_X$ y
$\mathcal G'=\mathcal O_{X_\eta}/\mathcal J'$; si
$\overline{\mathcal G'}=\mathcal O_X/\mathcal J$, se tiene, en efecto,
$\mathcal J'=i^*(\mathcal J)\mathcal O_{X_\eta}$, por tanto
$i^{-1}(\overline{Z'})=Z'$ (\textbf I, 4.4.5).
\end{proof}

Se observará que el preesquema $\overline{Z'}$ es la \emph{imagen cerrada} de $Z'$
por el morfismo compuesto
$Z'\longrightarrow X_\eta\xrightarrow{i}X$,
donde la primera flecha es la inyección canónica (\textbf I, 9.5.3); su espacio
subyacente es la clausura \emph{en $X$} de $Z'$ (\textbf I, 9.5.4), lo que
justifica la notación adoptada. Se dice también que $\overline{Z'}$ es el
subpreesquema de $X$ \emph{clausura} de $Z'$.
